%% ============================================================
%%  chapters/chapter2.tex - Literature Review and Related Work
%% ============================================================
\chapter{Literature Review and Related Work}
\label{chap:literature}

\section{Introduction}
\label{sec:lit-intro}

This chapter reviews the specialized knowledge areas required for the dissertation: keystroke dynamics, continuous authentication, deep learning sequence embeddings, programming trace analytics, online assessment integrity, and ethical deployment of behavioural biometrics. The chapter ends with a comparative analysis that establishes the knowledge gap addressed by BLAIM.

\section{Keystroke Dynamics as Behavioural Biometrics}
\label{sec:lit-keystroke}

Keystroke dynamics uses measurable typing behaviour for recognition. Common features include dwell time, flight time, press latency, release latency, key order, digraph latency, typing speed, and rhythm consistency. Unlike fingerprint or face biometrics, keystroke biometrics can be collected through an ordinary keyboard while a user performs normal work. This makes it suitable for online learning environments where additional hardware is expensive or intrusive.

Monrose and Rubin established that typing patterns can support biometric authentication and introduced important timing features for identity recognition \cite{monrose2000keystroke}. Later research expanded the evidence base by studying larger and more diverse datasets. Dhakal et al. analysed 136 million keystrokes and showed that typing behaviour contains both individual regularity and context-dependent variation \cite{dhakal2018observations}. This dual property is central to BLAIM: the system must use typing rhythm as identity evidence while acknowledging that stress, keyboard type, task type, and fatigue can shift the rhythm.

Traditional systems often used statistical summaries and distance measures. These approaches are attractive because they are interpretable and lightweight, but they struggle when the typed text or task changes. Programming assessments add another challenge because students alternate between normal typing, symbols, indentation, deletion, navigation, and pauses for thinking. Therefore, a system for programming assessment must represent both low-level timing and higher-level session behaviour.

\subsection{Feature Families in Keystroke Biometrics}
\label{subsec:feature-families}

Keystroke research commonly separates features into timing, sequence, and error-correction families. Timing features include hold latency and flight latency. Sequence features represent the order of keys, digraphs, trigraphs, and symbol-heavy transitions. Error-correction features include backspaces, deletes, retyping, and correction bursts. For natural-language typing, alphabetic digraphs may dominate. For programming, symbol transitions such as brackets, semicolons, indentation, and modifier keys are more important. This is one reason BLAIM includes key code as a normalized input feature rather than relying only on aggregate dwell and flight statistics.

Deep learning studies such as Altwaijry's keystroke authentication work and Li et al.'s multi-stream convolutional approach show the broader movement from manually selected statistics toward learned representations \cite{altwaijry2020deep,li2020keystroke}. However, learned representations still require careful interpretation in education. A classifier may separate users, but an instructor also needs to understand whether the session shows normal problem solving, struggle, or suspicious process behaviour.

\section{Continuous Authentication}
\label{sec:lit-continuous-auth}

One-time authentication checks whether a user is legitimate at login. Continuous authentication checks whether the active user still resembles the enrolled user throughout a session. This is important in online examinations because the risk is not limited to initial login. A session may involve proxy typing, shared control, device switching, or copied content after the session starts.

Keystroke authentication has been studied in online examination settings because it can operate without cameras and without interrupting the learner \cite{mattsson2020keystroke}. Senarath et al. re-evaluated keystroke dynamics for continuous authentication and identified the practical value of repeated checking, while also noting sensitivity to environment and device changes \cite{senarath2023reevaluating}. These findings support the BLAIM design decision to perform both direct verification and rolling monitoring, and to expose confidence and risk instead of treating every mismatch as misconduct.

Continuous authentication systems must balance false acceptance and false rejection. False acceptance allows an impostor through. False rejection challenges a genuine student and can create unfair stress. In educational deployment, both errors have human consequences. This is why BLAIM uses human-in-the-loop review and records evidence for interpretation rather than issuing automatic punishment.

\subsection{Verification and Identification}
\label{subsec:verification-identification}

Two biometric tasks are relevant to BLAIM. Verification is a 1:1 question: does this session match the claimed student? Identification is a 1:N question: which enrolled student does this session most resemble? Verification supports normal assessment login and monitoring. Identification supports investigation and audit workflows by showing the closest enrolled profiles for a session. BLAIM implements both because they answer different review questions and improve the depth of evidence available to instructors.

\section{Deep Learning and TypeNet}
\label{sec:lit-typenet}

Deep learning has improved keystroke biometrics by learning representations from event sequences rather than relying only on hand-crafted statistics. Long Short-Term Memory networks are well suited to this problem because typing is sequential and the meaning of a timing value depends on nearby events. Hochreiter and Schmidhuber introduced LSTM networks to address long-range sequence learning limitations in standard recurrent networks \cite{hochreiter1997lstm}.

TypeNet is a major reference point for this dissertation. Acien et al. proposed a deep learning keystroke biometric model that processes fixed-length keystroke sequences and produces embeddings for comparison \cite{acien2022typenet}. The embedding approach is useful because new users can be enrolled by storing a template rather than retraining a classifier for every class. This matches educational platforms where new students are added every semester.

BLAIM uses a TypeNet-inspired pipeline. Each keystroke sequence has 70 events and five features: hold latency, inter-key latency, press latency, release latency, and normalized key code. The model outputs a 128-dimensional embedding. Enrollment stores the mean and standard deviation of multiple embeddings per student and phase. Verification compares a new embedding against the stored template using cosine similarity.

\section{Programming Trace Analytics}
\label{sec:lit-trace}

Programming assessments produce rich process traces. A final code file can show what the student submitted, but the interaction trace can show how the solution was built. This is important for both learning support and integrity review.

Arakawa et al. studied in situ identification of self-regulated learning struggles in programming assignments and showed that trace behaviour can reveal meaningful struggle moments \cite{arakawa2021situ}. Dong et al. used student trace logs to identify meaningful progress and struggle during programming problem solving \cite{dong2021trace}. These studies support the idea that deletion bursts, long pauses, rewriting, and non-linear progress can be educationally meaningful.

BLAIM adopts this view by treating keystroke data as a process record. A high deletion rate may indicate confusion or iterative debugging. Long pauses may indicate planning, external help, distraction, or cognitive load. Paste events may indicate efficient reuse or suspicious bulk insertion depending on context. Therefore, BLAIM reports multiple indicators rather than reducing behaviour to a single misconduct label.

\subsection{From Integrity Detection to Learning Support}
\label{subsec:integrity-learning}

An important educational distinction is that struggle is not misconduct. A student who pauses, deletes, rewrites, and slowly improves may show a strong learning process. A student who submits correct code with almost no typing, no correction, and large paste-like insertions may require review even if the final code passes tests. BLAIM therefore frames behaviour analysis as a dual-purpose signal. It can support academic integrity triage, but it can also help instructors identify students who need formative support.

Recent work on LLM-assisted formative feedback shows that code context and interaction evidence can be transformed into useful feedback for learners \cite{nguyen2023gpt4}. BLAIM follows this direction by allowing optional LLM-assisted explanation while keeping rule-based metrics available when no LLM is configured.

\section{Academic Integrity in Online Assessment}
\label{sec:lit-integrity}

Remote assessment has increased reliance on proctoring, lockdown browsers, plagiarism detection, and identity checks. However, many proctoring methods raise privacy and equity concerns. Webcam proctoring can be intrusive, depends on stable network and private space, and may disadvantage students with limited resources or accessibility needs. Rapchak discusses remote proctoring technologies and the tension between monitoring and student privacy \cite{rapchak2020remote}.

Recent webcam-proctoring studies confirm that monitoring can reduce the opportunity or perceived opportunity for misconduct. Hylton et al. found that webcam-based proctoring reduced students' perceived opportunity to collaborate or use unauthorized resources, even though score differences between proctored and unproctored groups were not statistically significant in their setting \cite{hylton2016webcam}. Alguacil et al. later reported, through a randomized field experiment, that webcam-monitored students obtained significantly lower scores than unmonitored students under otherwise similar online exam conditions, interpreting the difference as evidence that monitoring can deter dishonest practices \cite{alguacil2023academic}. These findings are important because they show that visual monitoring has deterrent value.

However, the same studies also show why a keystroke-based model is needed for programming assessments. Webcam proctoring mainly observes the student and environment; it does not directly model how the submitted code was created. It can discourage collaboration and unauthorized resources, but it provides limited fine-grained evidence about authorship continuity, typing rhythm, copy-paste behaviour, pause structure, deletion patterns, or struggle and recovery during coding. It also depends on camera availability, lighting, network stability, private physical space, and student acceptance of video surveillance. Hylton et al. explicitly position webcam proctoring as a deterrent layer, while Alguacil et al. describe it as a low-touch surveillance intervention; neither approach produces the behavioural production trace that a programming instructor needs for process-aware review \cite{hylton2016webcam,alguacil2023academic}.

Keystroke dynamics is therefore a strong complement and, in some programming-assessment contexts, a more suitable primary signal. It is less visually intrusive than camera-based proctoring, works inside the normal coding environment, requires no additional biometric hardware, and continuously captures authorship evidence while the work is being produced. In programming assessments, it can complement code similarity tools by showing whether code appeared gradually through natural typing, through iterative debugging, or through sudden paste-like behaviour. This is the central reason BLAIM focuses on keystroke authentication and behaviour analysis rather than relying on webcam monitoring alone.

\subsection{Comparison with Static Plagiarism Detection}
\label{subsec:static-plagiarism}

Static code similarity tools examine the submitted artifact. They are strong when copied code is similar to known sources or peer submissions, but weaker when the external source is private, generated, heavily modified, or unavailable. Keystroke evidence examines the production process instead. It cannot prove semantic originality by itself, but it can show whether the final code was constructed through a plausible human editing process. The strongest integrity workflow combines both: static similarity for product evidence and keystroke dynamics for process evidence.

\section{Ethical, Security, and Fairness Considerations}
\label{sec:lit-ethics}

Keystroke data is sensitive because it can identify a person through behaviour. Responsible deployment requires informed consent, data minimization, access control, retention limits, and clear appeal processes. Biometric systems can also show uneven performance across populations or contexts. Research on fairness in biometric systems, such as the work by Buolamwini and Gebru, demonstrates that accuracy disparities can occur when systems are not evaluated across diverse groups \cite{buolamwini2018gender}. Although that study focused on face analysis, the fairness lesson applies to keystroke biometrics: institutions must test error rates across relevant learner groups before high-stakes use.

BLAIM therefore positions its outputs as decision-support evidence. It should not be used as an automatic disciplinary engine. A high-risk event should trigger review, context gathering, and student response, not immediate penalty.

\subsection{Security Requirements for Biometric Templates}
\label{subsec:security-templates}

Biometric templates are sensitive even when raw passwords or full text are not stored. A keystroke template can still represent a student's behavioural identity. Production systems should therefore apply access control, audit logging, encryption at rest, retention schedules, and administrative separation of duties. The BLAIM prototype stores templates and archives in PostgreSQL to demonstrate functionality, while Chapter 4 identifies encryption and access hardening as production requirements.

\section{Microservice Deployment Literature}
\label{sec:microservice-lit}

The implementation is also informed by microservice architecture principles. Newman describes microservices as independently deployable services organized around business capabilities \cite{newman2015microservices}. Balalaie et al. discuss how microservice architecture supports DevOps and continuous delivery \cite{balalaie2016microservices}. These ideas fit GradeLoop because authentication, assessment, evaluation, and analytics can evolve independently while communicating through APIs and events. BLAIM follows this model by implementing keystroke intelligence as a separate FastAPI service with its own database schema, Redis session buffer, and integration endpoints.

\section{Comparison of Related Work}
\label{sec:comparison}

Table~\ref{tab:related-comparison} compares selected related approaches against the requirements of this dissertation.

\begin{table}[H]
  \centering
  \caption{Comparison of Related Work and BLAIM}
  \label{tab:related-comparison}
  \begin{singlespace}
  \begin{tabularx}{\textwidth}{@{}p{3.0cm}p{2.4cm}p{2.5cm}p{2.7cm}p{3.0cm}@{}}
    \toprule
    \textbf{Work or system type} & \textbf{Continuous identity} & \textbf{Learning behaviour} & \textbf{Deployable LMS service} & \textbf{Key limitation addressed by BLAIM} \\
    \midrule
    Classical keystroke biometrics \cite{monrose2000keystroke} & Identity-focused & No & No & Adds education-specific behaviour analytics and service integration \\
    Large-scale typing analysis \cite{dhakal2018observations} & No & Typing-focused & No & Converts typing insights into an operational assessment module \\
    TypeNet \cite{acien2022typenet} & Yes & No & No & Uses embeddings inside an LMS-ready pipeline with dashboards and archives \\
    Continuous exam authentication \cite{senarath2023reevaluating} & Yes & Integrity-focused & Research-oriented & Adds process interpretation, multi-phase enrollment, and GradeLoop integration \\
    Programming trace analytics \cite{arakawa2021situ,dong2021trace} & No & Yes & Research-oriented & Adds biometric identity continuity to process analytics \\
    BLAIM implementation & Yes & Yes & Yes & Integrates biometric and behavioural evidence from one keystroke stream \\
    \bottomrule
  \end{tabularx}
  \end{singlespace}
\end{table}

\section{Identified Knowledge Gap}
\label{sec:lit-gap}

The reviewed literature shows strong work in separate areas, but a practical gap remains. Keystroke biometrics can authenticate users, but often lacks pedagogical explanation. Programming trace analytics can explain learning behaviour, but usually lacks identity continuity. Remote proctoring can provide oversight, but raises privacy and equity concerns. BLAIM addresses this gap by combining continuous keystroke authentication and behaviour analysis in one microservice architecture designed for programming assessment workflows.

\clearpage
